\newpage
\phantomsection
\label{prf:constant-multiple}

\begin{remark*}[Return]
\hyperref[thm:constant-multiple-rule]{Return to Theorem}
\end{remark*}

\begin{theorem*}[Constant Multiple Rule]
Let $I \subseteq \mathbb{R}$ be an open interval, let $c \in I$, let $f : I \to \mathbb{R}$ be differentiable at $c$, and let $\alpha \in \mathbb{R}$. Then $\alpha f$ is differentiable at $c$ and $(\alpha f)'(c) = \alpha f'(c)$.
\end{theorem*}

\begin{proof}
\textbf{Professional Standard Proof.}~\\
Set $h := \alpha f$. For $x \in I$ with $x \neq c$, the difference quotient becomes
\[
  \frac{h(x) - h(c)}{x - c} = \frac{\alpha f(x) - \alpha f(c)}{x - c} = \alpha \cdot \frac{f(x) - f(c)}{x - c}.
\]
By the Constant Multiple Rule for Limits,
\[
  \lim_{x \to c} \frac{h(x) - h(c)}{x - c} = \alpha \cdot \lim_{x \to c} \frac{f(x) - f(c)}{x - c} = \alpha f'(c),
\]
where the final equality follows from the differentiability of $f$ at $c$. Therefore $h'(c) = \alpha f'(c)$.
\end{proof}

\begin{proof}
\textbf{Detailed Learning Proof.}~\\

\textbf{Step 1.} Form the difference quotient of $h := \alpha f$. For $x \in I$ with $x \neq c$, substituting the definition of $h$ gives
\[
  \frac{h(x) - h(c)}{x - c} = \frac{\alpha f(x) - \alpha f(c)}{x - c} = \alpha \cdot \frac{f(x) - f(c)}{x - c}.
\]
The scalar $\alpha$ factors out by basic arithmetic properties of real numbers.

\textbf{Step 2.} Apply the limit properties to establish differentiability. By the Constant Multiple Rule for Limits, the scalar $\alpha$ can be moved outside the limit:
\[
  \lim_{x \to c} \frac{h(x) - h(c)}{x - c} = \alpha \cdot \lim_{x \to c} \frac{f(x) - f(c)}{x - c}.
\]
Since $f$ is differentiable at $c$, the limit $\lim_{x \to c} \frac{f(x) - f(c)}{x - c}$ exists and equals $f'(c)$ by definition of differentiability at a point. Therefore $h'(c) = \alpha f'(c)$.
\end{proof}

\begin{remark*}[Proof structure]
This is a direct proof using the definition of differentiability. The strategy exploits the linearity of limits: the scalar $\alpha$ factors out of the difference quotient algebraically, then the Constant Multiple Rule for Limits allows passage of the scalar through the limit operation. The differentiability of the original function $f$ supplies the required limit value.
\end{remark*}

\begin{remark*}[Dependencies]~\\
\begin{itemize}
  \item \hyperref[def:derivative-at-point]{Definition of Differentiability at a Point}
  \item \hyperref[thm:limit-constant-multiple]{Constant Multiple Rule for Limits}
\end{itemize}
\end{remark*}